\Subhead{2.3.1}{Оптические кристаллы. Уравнение Френеля. Обыкновенная и необыкновенная волны}

В анизотропном диэлектрике $D_i=\varepsilon_0\varepsilon_{ij}E_j$. Из уравнений Максвелла при $\mu_r=1$ получаем
\begin{equation}
 \left[n^2\delta_{ij}-n_in_j-\varepsilon_{ij}\right]E_j=0,
 \qquad \det\left[n^2I-\bm n\bm n^{\mathsf T}-\varepsilon\right]=0,
 \quad \bm n=c\bm k/\omega.
\end{equation}
Это волновая форма уравнения Френеля. Для заданного направления $\bm k$ оно определяет две электромагнитные нормальные волны. Поперечен $\bm D$, поскольку $\bm k\cdot\bm D=0$; $\bm E$ вообще не обязан быть поперечным.

В одноосном кристалле $\varepsilon=\operatorname{diag}(n_o^2,n_o^2,n_e^2)$, ось $z$ — оптическая. Обыкновенная волна имеет $k^2=n_o^2\omega^2/c^2$ и $\bm E$ вне главной плоскости $(\bm k,z)$. Необыкновенная удовлетворяет
\begin{equation}
 \frac{k_x^2+k_y^2}{n_e^2}+\frac{k_z^2}{n_o^2}=\frac{\omega^2}{c^2},\quad
 \frac1{n_{\rm eff}^2(\theta)}=\frac{\sin^2\theta}{n_e^2}+\frac{\cos^2\theta}{n_o^2}.
\end{equation}
При распространении по оптической оси обе скорости совпадают. Нормаль к изочастотной поверхности задаёт $\bm v_g$; для эллипсоида она обычно не параллельна радиусу $\bm k$. Отсюда снос необыкновенного пучка. Разложение падающего поля по двум собственным поляризациям объясняет двойное лучепреломление и фазовые пластинки.
